\paragraph{What worked.}
\begin{itemize}
\item \emph{FLAP reproduces on a small edge model.} On SmolLM2-1.7B the ranking FLAP $<$ Wanda-SP $<$ Magnitude-SP in perplexity holds at every ratio tested, as the FLAP paper reports for LLaMA models of 7\,B parameters and larger. Light pruning is cheap in quality: FLAP-10\% moves WikiText-2 perplexity only from 18.48 to 20.82.
\item \emph{Pruning gives measurable resource savings on the GPU.} At 50\%, energy per generated token falls by 27\% (2.09 to 1.52\,J), peak memory from 6.48 to 5.00\,GB, and prefill latency by 35\% (512 tokens) and 38\% (1024 tokens).
\item \emph{Per-query routing has real, if small, headroom.} After excluding every variant that does not beat a constant answer, the oracle still exceeds the best fixed variant by 1.9 points (64.1\% to 66.1\%) at the same cost, with gains on PIQA $+$5.2, ARC-Easy $+$4.6, GSM8K $+$4.2, HellaSwag $+$2.8. On PIQA, FLAP-50\% keeps 60.8\% accuracy (constant answer: 50.1\%) with 47\% fewer parameters.
\item \emph{A method that prevents an overclaim.} The plain oracle would suggest a routing gain of 11.7 points; most of it comes from variants that answer every question identically. The informative oracle introduced here measures the gain a router could actually exploit.
\item \emph{Reliable infrastructure.} Results are identical question by question across runs on different days; every result is cross-checked against lm-eval; several silent errors that would have corrupted results were found and fixed (Appendix~\ref{app:issues}).
\end{itemize}
\paragraph{What did not work.}
\begin{itemize}
\item Accuracy falls steeply: macro accuracy 62.5\% (unpruned) to 47.3\% (20\%), 40.0\% (35\%) and 32.6\% (50\%). The most frequent MMLU answer covers 100\% of questions for FLAP-35\% and 100\% for FLAP-50\% (unpruned: 32\%), and on GSM8K the pruned variants reach the 256-token limit on 80--100\% of questions.
\item On the T4 at batch size 1, decode latency stays at 30.4--31.1\,ms per token for every ratio, so pruning does not shorten generation there; with similar multiple-choice costs and longer GSM8K outputs, per-query routing saves no latency on this GPU.
\item Laptop-CPU timings were not reliable enough to report (Section~\ref{sec:results}).
\end{itemize}
\paragraph{What is next.} The results point to concrete steps: evaluating milder ratios (10--30\%) to test whether the pruned variants stay above the constant-answer baseline on more tasks; removing the per-token overhead on the GPU (static KV cache with CUDA graphs) and int8 weights on the CPU, both implemented and unit-tested but not yet measured; and a controlled CPU measurement. The remaining phases are planned as follows.
